% ============================================================
%  深圳大学实验报告模板（SZU Experiment Report Template）
% This template can be modified according to the specific requirements of your experiment or project.
% Produced by: Chao Fan and Kewei Ou, AI School of SZU
% ============================================================

\documentclass[a4paper,12pt]{article}

% ----------------------------
% 中文与字体支持
% ----------------------------
\usepackage{fontspec}    % 字体支持
\usepackage{xeCJK}       % 中文字体支持
\setCJKmainfont{Noto Serif CJK SC} % 设置中文字体（Overleaf自带Noto字体）

% ----------------------------
% 页面设置与常用宏包
% ----------------------------
\usepackage{geometry}    % 页面边距控制
\geometry{left=1in, right=1in, top=1in, bottom=1in}

\usepackage{longtable}   % 支持长表格
\usepackage{graphicx}    % 插入图片
\usepackage{fancyhdr}    % 页眉页脚控制
\usepackage{tikz}        % 绘制框线等图形
\usetikzlibrary{calc}    % 坐标计算
\usepackage{verbatim}    % 显示代码块
\usepackage{float}       % 控制图片浮动位置（[H]参数）
\usepackage{xcolor}      % 标红本人参与部分
\newcommand{\me}[1]{\textcolor{red}{#1}}   % 用法：\me{短语}

% ----------------------------
% 页眉页脚设置
% ----------------------------
% \pagestyle{fancy}
\fancyhf{} % 清空默认页眉页脚
% \fancyhead[L]{深圳大学实验报告} % 左侧页眉文字
\fancyhead[C]{} % 中间空
\fancyhead[R]{} % 右侧空

% ============================================================
%                      文档开始
% ============================================================

\begin{document}

% ============================================================
% 封面页
% ============================================================
\begin{titlepage}
    \centering
    \vspace*{2cm}
    \Huge{\textbf{深 \ 圳 \ 大 \ 学 \ 实 \ 验  \ 报 \ 告}}\\[1.5cm]
    
    {\Large\textbf{课程名称}\quad\underline{\makebox[8cm][c]{Python程序设计}}}\\[0.5cm]
{\Large
\noindent\makebox[3.2cm][l]{\textbf{项目名称}}
\underline{\makebox[8cm][c]{Python阶段性作业1 Numpy}}}\\[0.5cm]
{\Large
\noindent\makebox[3.2cm][l]{\textbf{学 \quad \quad 院}}
\underline{\makebox[8cm][c]{人工智能学院}}}\\[0.5cm]
{\Large
\noindent\makebox[3.2cm][l]{\textbf{专 \quad \quad 业}}
\underline{\makebox[8cm][c]{人工智能}}}\\[0.5cm]
{\Large
\noindent\makebox[3.2cm][l]{\textbf{指导教师}}
\underline{\makebox[8cm][c]{舒婷、樊超}}}\\[0.5cm]
{\Large
\noindent\makebox[3.2cm][l]{\textbf{报 \quad 告 \quad 人}}
\underline{\makebox[3cm][c]{黄淇典}}
\hspace{0.5cm}
\textbf{学号}
\underline{\makebox[3cm][c]{2025305011}}}\\[0.5cm]
{\Large
\noindent\makebox[3.2cm][l]{\textbf{实验时间}}
\underline{\makebox[8cm][c]{2026/11/15}}}\\[0.5cm]
{\Large
\noindent\makebox[3.2cm][l]{\textbf{提交时间}}
\underline{\makebox[8cm][c]{2026/11/21}}}\\[1.5cm]

    \vfill
    \Large{\textbf{教务处制}}
\end{titlepage}

\newpage

% ============================================================
% 正文部分
% ============================================================
% Experiment Objectives


\section{Objectives}

本次实验进行Numpy的学习，\me{实现}一个image—coder的python代码编写，将图片转化成数据方便后续识图处理。终极目标是学会使用Python，感受transformer架构的魅力，成为杰出的程序员


% Experiment Process
\section{Overview}
实验采用自底向上的方式：先在 image-encoder.py 中\me{补全六个 TODO}，依次实现图像分块、线性投影与 CLS 拼接、多头自注意力、Transformer block 堆叠和输出头；再借助 check.py 按八个阶段\me{逐级验证}。每一步都先推导输入输出 shape，再编写代码，最后\me{比对数值误差}，形成"先看形状、再写算术"的调试习惯，确保出错时能快速定位到具体环节。
\begin{enumerate}
    \item[Spot 1]: 输入与分块。图像形状为 [B,3,224,224]，按 32×32 切成 7×7 共 49 个 patch，每块展平为 3×32×32 共 3072 维。先把末两轴拆成格行、格内行、格列、格内列，得到六轴 [B,3,7,32,7,32]，再调整轴顺序为 [B,7,7,3,32,32]，最后合并末三轴得 [B,49,3072]。由于图像已按给定均值与方差归一化，这里只做形状变换，不改变数值；reshape 不改变元素总数，可用 size 守恒\me{自查};
    \item[Spot 2]: Transformer block。49 个 patch token 与 CLS 拼成 50 个 token，每 token 768 维，先做 LayerNorm，再按 12 个头拆成 [B,12,50,64]。注意力分数由 q 与 k 的最后两轴交换后相乘得到 [B,12,50,50]，除以 √64 缩放并经 softmax 得系数，乘 v 得上下文 [B,12,50,64]。合并 12 个头后经输出投影加残差，再经 GELU MLP 与第二次残差，重复 12 次。多头拆分让不同头关注不同特征子空间，所有头由同一投影产生;
    \item[Spot 3]: 输出头。取第 0 个 CLS 向量 [B,768]，经 post-norm 归一化后乘投影矩阵 visual.proj 得 [B,512]。此处与其他线性层不同：普通权重按 [out,in] 存储，相乘前必须转置；而 visual.proj 按 [in,out] 存储，直接相乘即可，转置会让 shape 与数值同时出错。最后对 512 维向量单位化，使每行模长为 1，便于后续做相似度比较。这一步决定了后续检索使用的图像表征质量。
\end{enumerate}

% Figure
\begin{figure}[t]
    \centering
    \includegraphics[width=0.92\linewidth]{figs/vit_overview.png}
    \caption{xxx}
    \label{fig:vit-arch}
\end{figure}

% Implementation Details
\section{Implementations}
代码改动集中在 image-encoder.py 的六个函数，权重加载、归一化、GELU、softmax 与单位化等辅助函数由课程提供，接口与 dtype 保持不变。图 1 左侧给出主干八步的完整流水线，右侧放大单个 block 内部八小步，以颜色区分本人实现的 TODO 与课程提供的部分。下面按分块、注意力、输出三个部分说明关键实现，每步都标注了对应的 shape 变化。As shown in Figure.~\ref{fig:vit-arch},  

\begin{enumerate}
    \item[Part 1]: patchify（TODO-01）。先取 images.shape 得到 B、C、H、W，令 p=32 并算出网格数 gh=gw=7，把图像 reshape 成 [B,3,7,32,7,32]，六轴含义依次为批、通道、格行、格内行、格列、格内列。接着用 transpose(0,2,4,1,3,5) \me{重排}成 [B,7,7,3,32,32]，让同一 patch 的通道与块内像素相邻，最后 reshape 成 [B,49,3072]。若把六轴顺序写反会打乱块内像素排列；全程 size 不变，是\me{检查轴顺序}是否正确的最快方法。;
    \item[Part 2]: 多头注意力（TODO-02/03）。分头 先把 [B,N,768] reshape 成 [B,N,12,64]，再 transpose(0,2,1,3) 把头轴提前得 [B,12,N,64]，12 个头即可并行计算。k 用 swapaxes(-1,-2) 只交换最后两轴，与 q 相乘得分数 [B,12,N,N]；\me{缩放写成} np.sqrt(np.float32(64)) 避免 float64 类型提升。softmax 沿最后一轴归一化，系数乘 v 得上下文，再转置合并回 [B,N,768]，最后经输出投影加偏置。注意 k 不能用 .T，否则批轴与头轴会被一并交换。;
    \item[Part 3]: block 与输出头（TODO-04/06）。block 残差先做 LayerNorm、注意力，输出投影并\me{加回原始} \(x\)，形成第二次残差；再经 LayerNorm、线性层、GELU、线性层加回第二次结果。补丁嵌入用 \texttt{patches} 乘投影矩阵且\me{不加偏置}，再拼 CLS 并\me{加位置嵌入}。输出时取 CLS 经 \texttt{post\_norm}，乘 \texttt{visual.proj}（不转置）得 \([B,8,512]\)，最后归一化。
\end{enumerate}


\section{Results}
用 check.py --stage all \me{逐级比对}，八个阶段全部通过。所有数值均为\me{本地实跑}结果，未使用参考输出作为输入；跟老师的模板\me{对比基本一致}，权重与图片随附件提供，刚刚的\me{插图}展示了完整 shape 流水线。As shown in Figure.~\ref{fig:vit-arch},  

\begin{figure}[H]
    \centering
    \includegraphics[width=0.92\linewidth]{figs/fig_pipeline.png}
    \caption{a beautiful picture of hqd's experimetn}
    \label{fig:vit-pipeline}
\end{figure}

\section{Discussion}
数值精度满足要求，残余误差来自 FP32 累加顺序与 softmax 细节，属正常范围。\me{基本没什么问题}，性能开销主要在 [B,H,N,N] 注意力矩阵与 12 个 block 顺序执行，后续可\me{尝试分块计算}降低内存占用。

\newpage

% ============================================================
% 批阅与成绩评定页
% ============================================================

% 绘制正文外框（包含批阅区域）
\begin{tikzpicture}[remember picture, overlay]
  \draw[thick] 
    ($(current page.north west)+(2cm,-3cm)$) 
    rectangle 
    ($(current page.south east)+(-2cm,3.5cm)$);
\end{tikzpicture}

\vspace{1cm}

% 批阅区
\noindent \textbf{指导教师批阅意见：}
\vspace{6cm}
\hfill

\vspace{1cm}

\noindent \textbf{成绩评定：}
\vspace{2cm}
\hfill

\vspace{1cm}

\noindent
\hfill
\textbf{指导教师签字：}\quad\quad\quad\quad\quad\\[0.5cm]

\hfill
\textbf{年}\quad\quad
\textbf{月}\quad\quad
\textbf{日}\quad

\vspace{1cm}

% 备注部分
\noindent \textbf{备注：}

% 框外的注释
\begin{tikzpicture}[remember picture, overlay]
  \node[
    anchor=north west,
    inner sep=0pt,
    text width=\dimexpr\paperwidth-4cm\relax
  ] at ($(current page.south west)+(2cm,3.1cm)$) {
    \begin{minipage}{\dimexpr\paperwidth-4cm\relax}
      \small
      \noindent
      \textbf{注：}\;1、报告内的项目或内容设置，可根据实际情况加以调整和补充。\\[0.2cm]
      \phantom{\textbf{注：}\;}2、教师批改学生实验报告时间应在学生提交实验报告时间后 10 日内。
    \end{minipage}
  };
\end{tikzpicture}

% ============================================================
%                      文档结束
% ============================================================

\end{document}
